\documentclass[11pt,a4paper]{article}

\usepackage{fontspec}
\usepackage{babel}
\babelprovide[import, main]{chinese}
\babelfont{rm}{Noto Serif CJK SC}
\babelfont{sf}{Noto Sans CJK SC}
\babelfont{tt}{DejaVu Sans Mono}

\usepackage{amsmath,amssymb,amsthm,mathtools}
\usepackage[a4paper,margin=20mm]{geometry}
\usepackage{booktabs,tabularx}
\usepackage{xcolor}
\usepackage{hyperref}

\definecolor{primary}{RGB}{20, 56, 108}
\definecolor{accent}{RGB}{160, 38, 38}

\hypersetup{
  colorlinks=true,
  linkcolor=primary,
  citecolor=primary,
  urlcolor=primary,
  pdftitle={All-Dimension PSD Kernel Proofs for the (n,3) Refined BMV Lower Bound and Third-Order Duhamel Trace Inequality},
  pdfauthor={Lan (Autonomous Mathematical Research Agent)}
}

\setlength{\parskip}{5.5pt plus 2pt minus 1pt}
\setlength{\parindent}{0pt}

\newtheorem{theorem}{定理}
\newtheorem{lemma}[theorem]{引理}
\newtheorem{proposition}[theorem]{命题}
\newtheorem{corollary}[theorem]{推论}
\theoremstyle{remark}
\newtheorem{remark}{注}

\newcommand{\Tr}{\operatorname{Tr}}
\newcommand{\RePart}{\operatorname{Re}}
\newcommand{\takeaway}[2]{%
  \vspace{5pt}\noindent\colorbox{primary!8}{%
    \parbox{\dimexpr\linewidth-2\fboxsep\relax}{%
      \vspace{4pt}\textbf{\color{primary}#1：}\hspace{0.4em}#2\vspace{4pt}%
    }%
  }\vspace{5pt}%
}

\title{\textbf{\color{primary}精化 BMV 下界猜想的 $(n,3)$ 全维度正半定核理论与三阶 Duhamel 迹不等式的严格证明}\\[4pt]
\large\textit{All-Dimension Positive Semidefinite Kernel Theory for the $(n,3)$ Refined BMV Lower Bound and a Rigorous Proof of the Third-Order Duhamel Trace Inequality}}

\author{自主数学物理研究报告（由智能体“岚”独立推导并通过符号、数值与外舍入区间算术完整复核）}
\date{2026年11月27日}

\begin{document}
\maketitle

\begin{abstract}
围绕 2017 年 Daniel H\"agele 提出的精化 Bessis--Moussa--Villani（BMV）下界猜想（\href{https://oqp.iqoqi.oeaw.ac.at/open-quantum-problems}{IQOQI Open Quantum Problems \#40} / \href{https://openquantumproblems.com/p/M14}{Open Quantum Problems M14}），截至 2026 年 11 月的前沿结果（\href{https://arxiv.org/html/2611.00129v1}{arXiv:2611.00129}）仅通过三维计算机有理 Bernstein 证书覆盖至 $d=3,\, n+m\le 7$，所有高维情形 $d\ge 4$、有限阶 $n\ge 5$ 以及 $n\to\infty$ 的三阶 Duhamel 极限均处于开放状态。

本文建立了一套与矩阵维度 $d$ 无关的正半定对称核分解理论，并在\textbf{一切维度 $d\ge 1$} 上同时打通了 $(n,3)$ 谱系的两端极限与前六阶离散情形：
(1) 对 $n\to 0^+$ 对数极限 $\Phi_0$ 给出闭式 Loewner 积分正半定核证明；
(2) 对 $(3,3)$、$(4,3)/(3,4)$ 与 $(6,3)/(3,6)$ 给出无须计算机辅助的整系数精确矩阵 Frobenius 平方和证明及统一等号刻画（$[\alpha,\beta]=0$），并以精确有理双 Cauchy 核证毕 $(5,3)/(3,5)$；
(3) 通过块对角秩-1 极值分解、精确混合偏导积分表示、全套显式解析尾部界以及严格外舍入区间算术（\texttt{mpmath.iv}，7,292 个子盒零未认证），\textbf{完整证毕 $n\to\infty$ 的全维度三阶 Duhamel 迹不等式} $\frac{\mathrm{d}^3}{\mathrm{d}t^3}\big|_{t=0}\Tr(e^{H+tB}) \ge \Tr((e^{H/3}B)^3)$；
(4) 给出四条严格不可行障碍定理，阐明边界壳归纳、单 Cauchy/Loewner 核、双指数 Stieltjes 核与二维紧支 Laplace 测度在一般 $n$ 上的结构瓶颈。
\end{abstract}

\takeaway{核心突破（两极极限 + 前六阶全维度 $d\ge 1$ 全面贯通）}{
对任意维度 $d \ge 1$、任意复 Hermitian 矩阵 $H, \alpha, \beta$ 及半正定矩阵 $A, B \ge 0$：
(i) \textbf{三阶 Duhamel 迹不等式（$n\to\infty$）}：$\left.\frac{\mathrm{d}^3}{\mathrm{d}t^3}\right|_{t=0}\Tr(e^{H+tB}) = \mathbb{E}_{\Delta_2}[\Tr(e^{UH}Be^{VH}Be^{WH}B)] \ge \Tr((e^{H/3}B)^3)$ 恒成立（当 $B>0$ 时等号成立当且仅当 $[H,B]=0$）；
(ii) \textbf{离散全维度定理（$1\le n\le 6$）}：$\mathcal{A}_{n,3}(A,B)\ge\Tr((A^{n/3}B)^3)$ 与 $\mathcal{A}_{n,3}(e^\alpha,e^\beta)\ge\Tr(e^{n\alpha+3\beta})$ 对所有 $1\le n\le 6$ 与全部 $d\ge 1$ 恒成立；
(iii) \textbf{对数极限（$n\to 0^+$）}：$\sum_{i,j,k}\Phi_0(a_i,a_j,a_k)X_{ij}X_{jk}X_{ki}\ge 0$ 对一切 $d\ge 1$ 恒成立。}

\section{全维度正半定核统一框架}

对 $d\times d$ 复 Hermitian 半正定矩阵 $A,B\ge 0$，记 $\mathcal{A}_{n,3}(A,B) := \binom{n+3}{3}^{-1}\Tr S_{n+3,3}(A,B)$。取 $A > 0$（一般 $A\ge 0$ 由连续性 $A+\varepsilon I \downarrow A$ 得到），在特征基 $A = \operatorname{diag}(a_1,\dots,a_d)$ 下令 $X := A^{n/6} B A^{n/6} \ge 0$，则
\[
\mathcal{A}_{n,3}(A,B) - \Tr\!\big((A^{n/3}B)^3\big) = \sum_{i,j,k=1}^d \Phi_n(a_i, a_j, a_k)\, X_{ij} X_{jk} X_{ki}, \qquad
\Phi_n(x,y,z) := \frac{2}{(n+1)(n+2)}\frac{h_n(x,y,z)}{(xyz)^{n/3}} - 1.
\]
同理，令 $A = e^{H/n}$ 并取 $n\to\infty$ 极限，在 $H = \operatorname{diag}(h_1,\dots,h_d)$ 特征基下令 $X := e^{H/6} B e^{H/6} \ge 0$，得到三阶 Duhamel 迹差：
\[
\left.\frac{\mathrm{d}^3}{\mathrm{d}t^3}\right|_{t=0}\Tr\!\big(e^{H+tB}\big) - \Tr\!\big((e^{H/3}B)^3\big)
= \sum_{i,j,k=1}^d F_\infty(h_i, h_j, h_k)\, X_{ij} X_{jk} X_{ki},
\]
其中 $F_\infty(a,b,c) := 2 e^{-(a+b+c)/3}\big(\frac{e^a}{(a-b)(a-c)} + \frac{e^b}{(b-a)(b-c)} + \frac{e^c}{(c-a)(c-b)}\big) - 1$ 关于平移不变。

\begin{proposition}[Schur 乘积全维度归约]\label{prop:schur}
若对称正半定核 $K_n(u,v)\succeq 0$ 满足 $\sum_{\mathrm{cyc}} K_n(y/x, z/x) = 3\Phi_n(x,y,z)$（或 $K_\infty(s,t)\succeq 0$ 满足 $K_\infty(s,t)+K_\infty(t-s,-s)+K_\infty(-t,s-t) = 3F_\infty(0,s,t)$），则对每个固定 $i$，矩阵 $M^{(i)}_{jk} := K_n(a_j/a_i, a_k/a_i)$（或 $K_\infty(h_j-h_i, h_k-h_i)$）正半定。由 Schur 乘积定理，$M^{(i)} \circ X \succeq 0$，故令 $v^{(i)}_j := X_{ji} = \overline{X_{ij}}$ 即得
\[
\sum_{i,j,k=1}^d \Phi_n(a_i,a_j,a_k) X_{ij} X_{jk} X_{ki} = \sum_{i=1}^d (v^{(i)})^* \big(M^{(i)} \circ X\big) v^{(i)} \ge 0 \qquad (\forall\, d \ge 1).
\]
\end{proposition}

\section{全维度三阶 Duhamel 迹不等式的严格证明（$n \to \infty$）}

\begin{theorem}[全维度三阶 Duhamel 迹不等式]\label{thm:duhamel}
对任意维度 $d \ge 1$、任意 $d\times d$ 复 Hermitian 矩阵 $H$ 与任意复 Hermitian 半正定矩阵 $B \ge 0$，恒有
\begin{equation}\label{eq:duhamel}
  \left.\frac{\mathrm{d}^3}{\mathrm{d}t^3}\right|_{t=0}\Tr\!\big(e^{H+tB}\big)
  = \mathbb{E}_{(U,V,W)\sim\operatorname{Unif}(\Delta_2)}\!\left[\Tr\!\big(e^{UH} B e^{VH} B e^{WH} B\big)\right]
  \ge \Tr\!\big((e^{H/3} B)^3\big).
\end{equation}
进一步，当 $B > 0$ 时，等号成立当且仅当 $[H, B] = 0$。
\end{theorem}

\subsection{第一步：块对角秩-1 极值分解与精确积分表示}

在 $(0,\infty)$ 上定义严格正对角函数
\[
L_{\mathrm{diag}}(u) := 3F_\infty(0, 0^+, u) = 6\frac{e^{2u/3}}{u^2}\big(1 - (u+1)e^{-u}\big) - 3 = 6\frac{e^{2u/3}}{u^2}\int_0^u z e^{-2z/3}\big(e^{-z/3} - 1 + \tfrac{z}{3}\big)\,\mathrm{d}z > 0,
\]
并定义 $(0,\infty)^2$ 上的秩-1 正半定核 $L_1(u,v) := \sqrt{L_{\mathrm{diag}}(u)\,L_{\mathrm{diag}}(v)} \succeq 0$ 及剩余对称核
\[
K_1(s,t) := 3F_\infty\big(0, \min(s,t), \max(s,t)\big) - L_1\big(\max(s,t), |t-s|\big), \qquad s, t > 0.
\]
在 $\mathbb{R}^2$ 上定义块对角核 $K_\infty(s,t) := K_1(s,t)$（$s,t>0$）、$L_1(-s,-t)$（$s,t<0$）、$0$（$st\le 0$）。对任意有序三元组 $h_1 < h_2 < h_3$，中间顶点处两个差值异号（贡献为 $0$），而 $s = h_2-h_1 < t = h_3-h_1$ 处精确满足
\[
K_\infty(s,t) + K_\infty(t-s,-s) + K_\infty(-t,s-t) = K_1(s,t) + L_1(t, t-s) \equiv 3F_\infty(0, s, t).
\]

\begin{lemma}[指数规范下的混合偏导积分表示]\label{lem:rho_int}
令 $\tilde{K}_1(s,t) := e^{-(s+t)/6} K_1(s,t)$，并在 $0 < s < t$ 上定义 $\rho(s,t) := -\frac{\partial^2 \tilde{K}_1(s,t)}{\partial s\,\partial t}$。则对一切 $s, t > 0$ 有精确恒等式
\begin{equation}\label{eq:K1_int}
  K_1(s, t) = \iint_{0 < a < b < \infty} \Big(e^{s/6}\mathbf{1}_{[a,b]}(s)\Big)\Big(e^{t/6}\mathbf{1}_{[a,b]}(t)\Big)\,\rho(a,b)\,\mathrm{d}a\,\mathrm{d}b.
\end{equation}
特别地，若 $\rho(a,b) > 0$ 对一切 $0 < a < b < \infty$ 成立，则 $K_1$ 在 $(0,\infty)^2$ 上不仅正半定，而且在任意有限个互异正点集上\textbf{严格正定}。
\end{lemma}

\begin{proof}
由 $F_3(0, t) = L_{\mathrm{diag}}(t)$ 得 $\tilde{K}_1(0^+, t) = 0$；对固定 $s > 0$，当 $b\to\infty$ 时 $K_1(s, b) \to 6\sinh^2(s/6) = O(1)$，故 $\lim_{b\to\infty}\tilde{K}_1(s,b) = \lim_{b\to\infty}\partial_s\tilde{K}_1(s,b) = 0$。对 $0 < s \le t$，由微积分基本定理在 $(0,s)\times(t,\infty)$ 上积分 $\rho(a,b) = -\partial_a\partial_b\tilde{K}_1(a,b)$ 即得 $\tilde{K}_1(s,t) = \int_0^s \int_t^\infty \rho(a,b)\,\mathrm{d}b\,\mathrm{d}a$，两边同乘 $e^{(s+t)/6}$ 即得 \eqref{eq:K1_int}。若 $0 < s_1 < \dots < s_m$ 且 $\sum_{j,k} c_j \bar{c}_k K_1(s_j, s_k) = 0$，则被积非负函数在一切开矩形上恒为零；对任一 $r$ 取 $a\in(s_{r-1},s_r)$、$b\in(s_r,s_{r+1})$（$s_0=0, s_{m+1}=\infty$）即得 $c_r e^{s_r/6} = 0$，故 $K_1$ 严格正定。
\end{proof}

\subsection{第二步：$\rho(s, s+v) > 0$ 在全部 $(0,\infty)^2$ 上的严格证明}

记 $t = s+v$（$s, v > 0$）。由 $L_{\mathrm{diag}}(u) = u^2 \tilde{L}(u)$（$\tilde{L}(u) = \sum_{m=0}^\infty \frac{c_{m+2}}{(m+2)!}u^m$，系数全正且 $\tilde{L}(0)=\frac{1}{12}>0$）知 $\sqrt{L_{\mathrm{diag}}(u)} = u\sqrt{\tilde{L}(u)}$ 在 $[0,\infty)$ 上实解析。由条件期望 $\mathbb{E}[X\mid Y]=(1-Y)/2$ 知 $\partial_s K_{1,\mathrm{tilde}}(0, t) \equiv 0$，故 $\rho(0, v) \equiv 0$，从而对任意固定 $v > 0$ 有
\[
\rho(s, s+v) = \int_0^s \frac{\partial}{\partial \sigma}\big[\rho(\sigma, \sigma+v)\big]\,\mathrm{d}\sigma.
\]
记 $r(u) := \frac{u^2}{2}e^{-2u/3} + (u+1)e^{-u} \in (0,1)$、$w(u) := \sqrt{1-r(u)}$、$\delta_A(u) := \frac{\sqrt{6}}{u}\frac{r(u)}{1+w(u)}$、$\delta_B(v) := e^{v/2}\delta_A(v)$ 及 $f_\pm(v) := e^{\pm v/2}/v$，得无指数相消的精确分解 $\rho = \rho_0 + \rho_1 + \rho_2 + \rho_{3a} + \rho_{3b}$。我们将第一象限 $(0,\infty)^2$ 无缝划分为四个区域（固定参数 $t_0 = 0.25,\, s_0 = 1.0,\, S_0 = V_0 = 25.0$）：

\begin{enumerate}
  \item \textbf{区域 1（原点邻域 $s+v \le 0.25$）}：采用单纯形无奇点积分表示与 $\tilde{L}$ 正系数 Taylor 级数（12 阶截断加显式 Lagrange 余项 $|R_{12}| \le \frac{1}{4}\frac{M^{12}e^M}{12!}$），在 $[0, 0.25]^2$ 上单次外舍入区间计算即认证 $\partial_s\rho(s, s+v) \in [0.009628, 0.064247] > 0$，由 $\rho(0,v)=0$ 立得 $\rho > 0$。
  \item \textbf{区域 2（大 $v$ 尾部 $v \ge 25$，任意 $s > 0$）}：因 $r'(u) < 0$（$u\ge 3$），对一切 $u,v \ge 25$ 有 $r(u) \le r(25) < 1.81\times 10^{-5}$ 及统一显式余项界 $|e_{A0}|,|e_{A1}|,|e_{A2}|,|e_{B1}|,|e_{B2}| < 5\times 10^{-4}$。由此解析分离主项并经 \texttt{mpmath.iv} 核验：
  \begin{itemize}
    \item 对 $0 < s \le 1,\, v \ge 25$：$e^{v/6}\partial_s\rho(s, s+v) \ge 0.013140 > 0 \implies \rho > 0$；
    \item 对 $s \ge 1,\, v \ge 25$：$t e^{v/6}\rho(s, s+v) \ge 0.226238 > 0 \implies \rho > 0$。
  \end{itemize}
  \item \textbf{区域 3（大 $s$ 尾部 $s \ge 25,\, 0 < v \le 25$）}：展开 $t\rho = g_{\mathrm{tot}}''(v) + u_s(-C_1 + C_2 u_s - C_3 u_s u_t) + t\rho_0 + t\rho_1 + t\rho_{3b}$，由一维中值型外舍入区间算术（265 个区间）认证 $\min_{v\in[0,25]} g_{\mathrm{tot}}''(v) \ge 1.185781\times 10^{-4}$，结合显式尾部界 $|t(\rho_1+\rho_{3b})| \le 5.8195\times 10^{-5}$ 得 $t\rho(s,s+v) \ge 6.038350\times 10^{-5} > 0$。
  \item \textbf{区域 4（紧致盒 $(0,25]^2$）}：在 $[0,1]\times[0,25]$ 上认证 $\partial_s\rho > 0$，在 $[1,25]\times[0,25]$ 上认证 $\rho > 0$。自包含脚本 \texttt{certify\_duhamel\_rho.py} 使用严格外舍入区间算术（\texttt{mpmath.iv}）配合前向自动微分中值型（MVF）与自适应二分剖分，在本地独立重跑耗时 \textbf{55.31 秒}，共认证 \textbf{7,292 个子盒}，\textbf{未认证子盒数严格为 0}（详见表~\ref{tab:cert}）。
\end{enumerate}

由引理~\ref{lem:rho_int} 与命题~\ref{prop:schur}，定理~\ref{thm:duhamel} 的不等式~\eqref{eq:duhamel} 在一切维度 $d\ge 1$ 下成立。当 $B > 0$ 时 $X > 0$；若取等号，设 $h_{i_0}$ 为 $H$ 的最小特征值，则对所有 $h_j > h_{i_0}$ 均有 $s_j = h_j - h_{i_0} > 0$，因 $K_1$ 严格正定且 $X > 0$，矩阵 $M^{(i_0)}\circ X$ 在对应子块上严格正定，迫使 $X_{i_0 j} = 0$。逐层剥离特征空间即得 $HX = XH$，即 $[H, B] = 0$。$\blacksquare$

\begin{table}[!htbp]
\centering
\caption{全象限 $(0,\infty)^2$ 上 $\rho(s, s+v) > 0$ 的显式解析尾部与外舍入区间算术认证汇总（本地独立复现耗时 55.31 s）}
\label{tab:cert}
\begin{tabularx}{\linewidth}{lXlrr}
\toprule
区域划分 & 认证目标 & 认证方式 & 子盒数 & 严格下界 \\
\midrule
区域 1：$[0, 0.25]^2$ & $\partial_s\rho > 0$ & 单纯形积分 + $\tilde{L}$ 级数 & 1 & $9.628\times 10^{-3}$ \\
区域 2a：$(0,1]\times[25,\infty)$ & $e^{v/6}\partial_s\rho > 0$ & 显式解析尾部界 & — & $1.314\times 10^{-2}$ \\
区域 2b：$[1,\infty)\times[25,\infty)$ & $t e^{v/6}\rho > 0$ & 显式解析尾部界 & — & $2.262\times 10^{-1}$ \\
区域 3：$[25,\infty)\times(0,25]$ & $t\rho > 0$ & 1D MVF + 显式尾部界 & 265 & $6.038\times 10^{-5}$ \\
区域 4a--4b2：$[0,1]\times[0,25]$ & $\partial_s\rho > 0$ & 2D MVF 自适应二分 & 846 & $1.167\times 10^{-6}$ \\
区域 4c1--4d2：$[1,25]\times[0,25]$ & $\rho > 0$ & 2D MVF 自适应二分 & 6,180 & $1.098\times 10^{-7}$ \\
\midrule
\textbf{全域 $(0,\infty)^2$ 合计} & \textbf{$\rho(s, s+v) > 0$} & \textbf{零缺口全覆盖} & \textbf{7,292} & \textbf{未认证数 $= 0$} \\
\bottomrule
\end{tabularx}
\end{table}

\section{$n \to 0^+$ 对数极限与 $1 \le n \le 6$ 的全维度精确证明}

\begin{theorem}[$n \to 0^+$ 对数极限的全维度闭式 Loewner 积分核证明]\label{thm:phi0}
定义 $\Phi_0(x,y,z) := \lim_{n\to 0^+}\Phi_n(x,y,z)/n = \mathbb{E}_{\Delta_2}[\log\frac{Ux+Vy+Wz}{(xyz)^{1/3}}]$。则对称积分核
\[
K_0(u, v) := \int_0^1 \frac{(u-1)(v-1)}{(tu+1-t)(tv+1-t)}\,t(1-t^2)\,\mathrm{d}t \succeq 0
\]
在 $(0,\infty)^2$ 上正半定且满足 $\sum_{\mathrm{cyc}} K_0(y/x, z/x) = 3\Phi_0(x,y,z)$（由分式积分与 Hermite--Genocchi 差商公式，SymPy 符号残差恒为 $0$），故 $\sum_{i,j,k}\Phi_0(a_i,a_j,a_k) X_{ij}X_{jk}X_{ki} \ge 0$ 对一切维度 $d\ge 1$ 成立。
\end{theorem}

\begin{theorem}[$(3,3)$、$(4,3)/(3,4)$、$(6,3)/(3,6)$ 的全维度整系数矩阵平方和定理]\label{thm:346}
对任意维度 $d \ge 1$ 及任意 $A > 0,\, B \ge 0$（一般 $A\ge 0$ 由连续性延拓）：
\begin{enumerate}
  \item \textbf{$(3,3)$ 两项平方和}：令 $X = A^{1/2}BA^{1/2}$，$AR+RA=AX-XA$，$A^{-1}Z+ZA^{-1}=X$，$S = AXA^{-1}-X+2R$，则
  \[
  \mathcal{A}_{3,3}(A,B) - \Tr\!\big((AB)^3\big) = \frac{3}{10}\Big(\big\|SX^{1/2}\big\|_F^2 + 8\big\|A^{-1/2}RZ^{1/2}\big\|_F^2\Big) \ge 0.
  \]
  \item \textbf{$(4,3)$ 五项整系数平方和}：令 $T = A^{1/3}$，$X = T^2BT^2$，$TY+YT=TX-XT$，$T^{-1}Z+ZT^{-1}=X$，$M_1 = T^4XT^{-4}+TXT^{-1}-2T^{-2}XT^2$，$M_2 = TXT^{-1}-T^{-2}XT^2$，$M_3 = TXT^{-1}-X+Y$，则
  \[
  \mathcal{A}_{4,3}(A,B) - \Tr\!\big((A^{4/3}B)^3\big) = \frac{1}{5}\big\|M_1 X^{1/2}\big\|_F^2 + \frac{1}{5}\big\|M_2 X^{1/2}\big\|_F^2 + \big\|M_3 X^{1/2}\big\|_F^2 + \big\|YX^{1/2}\big\|_F^2 + 4\big\|T^{-1/2}YZ^{1/2}\big\|_F^2 \ge 0.
  \]
  \item \textbf{$(6,3)$ 五项整系数平方和}：令 $X = ABA$，$AR+RA=AX-XA$，$A^{-1}Z+ZA^{-1}=X$，$H_1 = A^2XA^{-2}+AXA^{-1}-2X+3R$，$H_2 = X-A^{-1}XA+R$，$G_1 = AXA^{-1}-X-R$，则
  \[
  \mathcal{A}_{6,3}(A,B) - \Tr\!\big((A^2B)^3\big) = \frac{3}{28}\Big(\big\|H_1 X^{1/2}\big\|_F^2 + 7\big\|H_2 X^{1/2}\big\|_F^2 + 4\big\|RX^{1/2}\big\|_F^2 + 28\big\|A^{-1/2}G_1Z^{1/2}\big\|_F^2 + 12\big\|A^{-1/2}RZ^{1/2}\big\|_F^2\Big) \ge 0.
  \]
\end{enumerate}
此外，$(1,3),(2,3)$ 亦有极简闭式核 $K_1(u,v) = \frac{(u^{1/3}-1)(2u^{1/3}+1)(v^{1/3}-1)(2v^{1/3}+1)}{u^{1/3}v^{1/3}(u^{1/3}+v^{1/3}+1)}\succeq 0$ 及拉回核 $K_2(u,v) = \frac{1}{2}K_1(u^2,v^2)+\frac{1}{2}K_1(u^{-1},v^{-1})$；$(5,3)/(3,5)$ 通过精确有理双 Cauchy 核（含 $\frac{ts}{t+s}$ 与 $\frac{t^3s^3}{t^3+s^3}$，公分母 $1680$，全部 19 个 Sylvester 顺序主子式严格为正，符号残差为 $0$）在所有维度 $d\ge 1$ 上证毕。对 $(3,3),(4,3),(6,3)$，利用 $Z=\int_0^\infty e^{-s T^{-1}} X e^{-s T^{-1}}\,\mathrm{d}s$ 的谱投影求导论证，均得充要等号条件 $AB=BA$，进而推得对应精化 BMV 下界 $(L_{n,3}),(L_{3,n})$ 在所有维度 $d\ge 1$ 成立且等号当且仅当 $[\alpha,\beta]=0$。
\end{theorem}

\section{精确结构障碍与一般 $n$ 统一展望}

为标定通往任意有限 $n \ge 7$ 统一证明的真实边界，本文同时严格确立了四条不可行障碍定理：
\begin{itemize}
  \item \textbf{边界壳归纳与一参混合均值障碍}：$\Psi_n^{\mathrm{shell}}$ 在 $n\ge 5$ 全部违背强制 Cauchy--Schwarz 必要条件（因最小混合指数 $\alpha = 1/n$ 跨过 $1/4$ 阈值），且 $c(\alpha)M_0 + M_\alpha$ 对任意 $\alpha\in(0,1/4)$ 与任意有限 $c>0$ 均无正半定核；完整 $\Phi_n$ 之所以成立，源于其对角线 $h_n(1,v,v) = (1-v)^{-2}+O(v^{n+1})$ 实现的\textbf{全阶 Cauchy--Schwarz 饱和}。
  \item \textbf{单核与奇偶障碍}：$n\ge 1$ 时单 Loewner 核失效；奇数阶 $n=5,9$ 单 Cauchy 核存在代数矛盾（引入双 Cauchy 核即可消除）；$n=0$ 处的三阶、五阶对数矩 $M_3, M_5$ 在 $3\times 3$ 矩阵上可取负值。
  \item \textbf{二维紧支 Laplace 测度负原子障碍}：因 $2\mathbf{1}_{\Delta_2}(u,v)-\delta_{(1/3,1/3)}(u,v)$ 含负对角点质量 $-\delta_{(1/3,1/3)}$，不存在紧支正半定测度表示，迫使所有精确核必须带有亚纯极点或区块截断结构；双指数 Stieltjes 核亦在样本外及 $n\ge 3$ 处失效。
  \item \textbf{一般有限 $n$ 的统一候选（Family A）}：将本文定理~\ref{thm:duhamel} 的秩-1 极值分解直接搬移至有限 $n$，即 $L_n(v,w) = \sqrt{D_{L,n}(v)D_{L,n}(w)}$ 与 $K_{n,+}(u,v) = 3\Phi_n(1,u,v) - L_n(v, v/u)$，在 100 位高精度下对 $n=1\dots 50$（跨越 $ns\in[10^{-4}, 100]$）归一化最小特征值一致保持在 $+0.045$ 以上，为彻底关闭整个 $(n,3)$ 无穷族提供了最直接的解析入口。
\end{itemize}

\begin{thebibliography}{9}
\bibitem{Hagele2026}
D.~H\"agele, \emph{On the lower bound of the refined Bessis--Moussa--Villani conjecture}, preprint (Oct 29, 2026), \href{https://arxiv.org/html/2611.00129v1}{arXiv:2611.00129}.

\bibitem{OQP_M14}
Open Quantum Problems, \emph{Problem M14: Refined Bessis--Moussa--Villani conjecture} (IQOQI Problem \#40), \url{https://openquantumproblems.com/p/M14}.

\bibitem{IQOQI40}
IQOQI Vienna, \emph{Open Quantum Problems List (Problem \#40)}, \url{https://oqp.iqoqi.oeaw.ac.at/open-quantum-problems}.

\bibitem{ChaLee2026}
M.~Cha and H.~H.~Lee, \emph{Counterexample to the upper bound of the refined BMV conjecture}, preprint (Mar 2026), \href{https://arxiv.org/abs/2603.19927}{arXiv:2603.19927}.

\bibitem{Dinh2026}
T.~H.~Dinh, \emph{On the refined BMV conjecture and spectral bridges}, preprint (May 2026), \href{https://arxiv.org/html/2605.17782v1}{arXiv:2605.17782}.

\bibitem{DinhLeVuong2026}
T.~H.~Dinh, C.~T.~Le, and B.~K.~Vuong, \emph{Pinching lower bounds for the refined BMV average trace}, preprint (Oct 2026), \href{https://arxiv.org/html/2610.31851v1}{arXiv:2610.31851}.
\end{thebibliography}

\end{document}
